% без водяного знака й без хендлу каналу
\def\nmtnowatermark{}
\documentclass[12pt]{article}
\usepackage[ukrainian,shorthands=off]{babel}   % shorthands=off: символ " не активний (потрібно для міток "..." у TikZ всередині \zadtask)
\usepackage{fontspec}
% --- НАЛАШТУВАННЯ ШРИФТІВ ---
% Шрифти Schoolbook-*.otf лежать у корені проєкту. Шлях підбирається автоматично:
% ../ (компіляція з папки теми), ./ (корінь проєкту), ../../ (підпапки теми 28); інакше -- TeX Gyre Schola.
\newcommand{\nmtSchoolbook}[1]{\setmainfont{Schoolbook-Regular.otf}[
    Path = #1,
    BoldFont = Schoolbook-Bold.otf,
    ItalicFont = Schoolbook-Italic.otf,
    BoldItalicFont = Schoolbook-BoldItalic.otf
]}
\IfFileExists{../Schoolbook-Regular.otf}{\nmtSchoolbook{../}}{%
 \IfFileExists{./Schoolbook-Regular.otf}{\nmtSchoolbook{./}}{%
  \IfFileExists{../../Schoolbook-Regular.otf}{\nmtSchoolbook{../../}}{%
   \setmainfont{texgyreschola}[Extension=.otf, UprightFont=*-regular, BoldFont=*-bold, ItalicFont=*-italic, BoldItalicFont=*-bolditalic]}}}
\usepackage[italic]{mathastext}
\usepackage[a4paper,margin=1.8cm,bottom=2cm,top=2cm]{geometry}
\usepackage{amsmath,amssymb}
\usepackage{tikz}
\usepackage{pgfplots}
\pgfplotsset{compat=1.18}
\usepgfplotslibrary{fillbetween}
\usetikzlibrary{calc,patterns,angles,quotes,intersections,babel,3d,shapes.symbols,shapes.geometric,decorations.pathreplacing,shadings,arrows.meta,hobby}
\usepackage{xcolor,array,fancyhdr,enumitem,multicol}
\usepackage{colortbl,multirow,diagbox,needspace}
\usepackage{tcolorbox}
\tcbuselibrary{skins,breakable}
\usepackage[hidelinks]{hyperref}
% водяний знак; щоб зібрати без нього (версія для вчителів), перед \documentclass пишуть \def\nmtnowatermark{}
\ifdefined\nmtnowatermark\else
\usepackage[firstpage=false, text=@pvtr2525, scale=0.6, color=gray!12]{draftwatermark}
\fi

% --- АВТОСТИСКАННЯ: рисунок/сітка, ширша за колонку, масштабується до ширини колонки ---
\newsavebox{\nmtfitbox}
\newenvironment{nmtfit}{\begin{lrbox}{\nmtfitbox}}{\end{lrbox}%
  \ifdim\wd\nmtfitbox>\linewidth \resizebox{\linewidth}{!}{\usebox{\nmtfitbox}}\else\usebox{\nmtfitbox}\fi}
\let\nmtIncludegraphicsOrig\includegraphics
\renewcommand{\includegraphics}[2][]{\begin{nmtfit}\nmtIncludegraphicsOrig[#1]{#2}\end{nmtfit}}


\definecolor{mainGreen}{RGB}{34, 120, 64}
\definecolor{yearOrange}{RGB}{220, 100, 30}
\definecolor{yearcolor}{RGB}{220, 100, 30}
\definecolor{titleBg}{RGB}{225, 240, 220}
\definecolor{instrBg}{RGB}{255, 240, 220}
\definecolor{instrBorder}{RGB}{220, 150, 80}
\definecolor{headerblue}{RGB}{0, 102, 204}   % використовується в рисунках старої бази

\setlength{\parindent}{0pt}
\emergencystretch=1.5em

\pagestyle{fancy}
\fancyhf{}
\renewcommand{\headrulewidth}{0.4pt}
\renewcommand{\footrulewidth}{0.4pt}
\fancyhead[C]{\small\color{gray!80} Пробний НМТ № 1}
\fancyhead[R]{}
\fancyfoot[L]{}
\fancyfoot[C]{\small\color{gray!80}\thepage}
\fancyfoot[R]{\small\color{gray!80}НМТ з математики}

% --- ТАБЛИЦІ ВІДПОВІДЕЙ ---
% ширина клітинки = (ширина рядка - відступи) / 5: на всю сторінку це ~3 см, у minipage поруч із рисунком таблиця стискається;
% п'ять колонок |m{w}| мають 10 відступів \tabcolsep і 6 вертикальних ліній -- тоді таблиця рівно на \linewidth
\newlength{\nmtcellw}
\newcommand{\nmtSetCell}{\setlength{\nmtcellw}{\dimexpr(\linewidth-10\tabcolsep-6\arrayrulewidth)/5\relax}}
\newcommand{\answerTable}[5]{%
\par\nopagebreak\vspace{0.15cm}
\noindent\nmtSetCell
\begin{tabular}{|*{5}{>{\centering\arraybackslash}m{\nmtcellw}|}}
\hline
\rule[-0.15cm]{0pt}{0.6cm}\textbf{А} & \rule[-0.15cm]{0pt}{0.6cm}\textbf{Б} & \rule[-0.15cm]{0pt}{0.6cm}\textbf{В} & \rule[-0.15cm]{0pt}{0.6cm}\textbf{Г} & \rule[-0.15cm]{0pt}{0.6cm}\textbf{Д} \\
\hline
\rule[-0.2cm]{0pt}{0.75cm}#1 & \rule[-0.2cm]{0pt}{0.75cm}#2 & \rule[-0.2cm]{0pt}{0.75cm}#3 & \rule[-0.2cm]{0pt}{0.75cm}#4 & \rule[-0.2cm]{0pt}{0.75cm}#5 \\
\hline
\end{tabular}
\par\vspace{0.25cm}
}
\newcommand{\answerTableTall}[5]{%
\par\nopagebreak\vspace{0.15cm}
\noindent\nmtSetCell
\begin{tabular}{|*{5}{>{\centering\arraybackslash}m{\nmtcellw}|}}
\hline
\rule[-0.2cm]{0pt}{0.7cm}\textbf{А} & \rule[-0.2cm]{0pt}{0.7cm}\textbf{Б} & \rule[-0.2cm]{0pt}{0.7cm}\textbf{В} & \rule[-0.2cm]{0pt}{0.7cm}\textbf{Г} & \rule[-0.2cm]{0pt}{0.7cm}\textbf{Д} \\
\hline
\rule[-0.45cm]{0pt}{1.2cm}#1 & \rule[-0.45cm]{0pt}{1.2cm}#2 & \rule[-0.45cm]{0pt}{1.2cm}#3 & \rule[-0.45cm]{0pt}{1.2cm}#4 & \rule[-0.45cm]{0pt}{1.2cm}#5 \\
\hline
\end{tabular}
\par\vspace{0.25cm}
}
\newcommand{\instructionBox}[1]{%
\par\vspace{0.3cm}
\begin{tcolorbox}[colback=instrBg, colframe=instrBorder, boxrule=0.6pt, arc=2pt,
    left=8pt, right=8pt, top=4pt, bottom=4pt]
\centering\small\bfseries #1
\end{tcolorbox}
\par\vspace{0.15cm}
}
\newcommand{\sectionTitle}[1]{%
\par\vspace{0.4cm}
\begin{tcolorbox}[colback=titleBg, colframe=titleBg, boxrule=0pt, arc=8pt,
    halign=center, fontupper=\Large\bfseries\color{mainGreen}]
\hyphenpenalty=10000 \exhyphenpenalty=10000 #1
\end{tcolorbox}
\par\vspace{0.3cm}
}
\newcommand{\task}[2]{\noindent\textbf{#1.}\ \ #2\par\vspace{0.2cm}}
% --- НМТ-СТИЛЬ ПОЛЕ ВІДПОВІДІ ---
\newcommand{\nmtAnswerBox}{%
\par\nopagebreak\vspace{0.2cm}\noindent
Відповідь:\ \
\begin{tabular}{@{}*{4}{|>{\centering\arraybackslash}p{0.8cm}}|@{\hspace{0.1cm}\textbf{,}\hspace{0.1cm}}*{3}{|>{\centering\arraybackslash}p{0.8cm}}|@{}}
\hline
\rule[-0.2cm]{0pt}{0.7cm} & & & & & & \\
\hline
\end{tabular}
\par\vspace{0.3cm}
}
% --- СІТКА ДЛЯ МАТЧИНГУ ---
\newsavebox{\nmtgridbox}
\newcommand{\matchingGrid}{%
    \sbox{\nmtgridbox}{\matchingGridRaw}%
    \ifdim\wd\nmtgridbox>\linewidth \resizebox{\linewidth}{!}{\usebox{\nmtgridbox}}\else\usebox{\nmtgridbox}\fi}
\newcommand{\matchingGridRaw}{%
    \begingroup
    \renewcommand{\arraystretch}{1.15}
    \setlength{\tabcolsep}{4pt}
    \footnotesize
    \begin{tabular}{r|c|c|c|c|c|}
         \multicolumn{1}{c}{} & \multicolumn{1}{c}{\textbf{А}} & \multicolumn{1}{c}{\textbf{Б}} & \multicolumn{1}{c}{\textbf{В}} & \multicolumn{1}{c}{\textbf{Г}} & \multicolumn{1}{c}{\textbf{Д}} \\ \cline{2-6}
         \textbf{1} & \phantom{X} & \phantom{X} & \phantom{X} & \phantom{X} & \phantom{X} \\ \cline{2-6}
         \textbf{2} & & & & & \\ \cline{2-6}
         \textbf{3} & & & & & \\ \cline{2-6}
    \end{tabular}
    \endgroup
}
% --- ДОДАТКОВІ МАКРОСИ ЗБІРНИКА ---
\newcommand{\taskBlock}[1]{%
\par\vspace{0.45cm}
\noindent{\color{mainGreen}\rule{\linewidth}{0.8pt}}\par\vspace{0.12cm}
\noindent{\small\bfseries\color{mainGreen}#1}\par\vspace{0.25cm}
}
\newcounter{zad}
\newcommand{\zadtask}[1]{\stepcounter{zad}\noindent\textbf{\thezad.}\ \ #1\par\nopagebreak\vspace{0.2cm}}
\newcommand{\zadnum}{\stepcounter{zad}\textbf{\thezad.}\ \ }
\newcounter{ans}
\newcommand{\ansitem}[1]{\stepcounter{ans}\noindent\textbf{\theans.}\ #1\par\vspace{0.07cm}}
% мітка року: нерозривна, притиснута праворуч; якщо не вміщається -- праворуч на новому рядку
\newcommand{\nmtyear}[1]{\unskip\nobreak\finalhyphendemerits=0 \hfill\penalty100\hskip1em\hbox{}\nobreak\hfill{\small\color{yearOrange}\mbox{\rule{0pt}{2.3ex}(НМТ~#1)}}}
% --- МАКРОС КОРОТКОГО РОЗВ'ЯЗКУ ---
\newcommand{\solution}[1]{%
\par\vspace{0.12cm}
\begin{tcolorbox}[colback=titleBg!45, colframe=mainGreen!55, boxrule=0.4pt, arc=2pt,
    left=6pt, right=6pt, top=3pt, bottom=3pt, breakable]
\footnotesize\textbf{Короткий розв'язок.}\ #1
\end{tcolorbox}
\par\vspace{0.12cm}
}
% Розв'язки й відповіді (є до завдань НМТ-2026): \showsolutionstrue -- показати, \showsolutionsfalse -- сховати
\newif\ifshowsolutions
\showsolutionsfalse
% --- ЗАГОЛОВКИ З ПУНКТАМИ КЛІКАБЕЛЬНОГО ЗМІСТУ ---
\setcounter{tocdepth}{2}
\newcommand{\chapterTitle}[1]{%
\clearpage\phantomsection\addcontentsline{toc}{section}{#1}%
\sectionTitle{#1}}
\newcommand{\typeTitle}[1]{%
\needspace{6\baselineskip}%
\phantomsection\addcontentsline{toc}{subsection}{\quad #1}%
\par\vspace{0.3cm}
\noindent{\large\bfseries\color{yearOrange}#1}\par\vspace{0.08cm}
\noindent{\color{yearOrange}\rule{\linewidth}{0.4pt}}\par\nopagebreak\vspace{0.2cm}}
\raggedbottom
% усередині одного завдання сторінка не розривається: нерозривний вертикальний проміжок
% і нерозривна межа абзаців (умова ніколи не відділяється від варіантів, рисунка чи сітки)
\newcommand{\nbvspace}[1]{\nopagebreak\vspace{#1}}
\newcommand{\nmtnobreak}{\ifvmode\penalty10000 \fi}
% --- ЄДИНИЙ МАКЕТ ЗАВДАНЬ НА ВІДПОВІДНІСТЬ ---
% мітка в колонці сталої ширини, текст -- у parbox: продовження довгого рядка
% вирівнюється під текстом, а не під міткою; відступ між пунктами однаковий скрізь
\newlength{\nmtMatchLab}\setlength{\nmtMatchLab}{1.6em}
\newlength{\nmtMatchSep}\setlength{\nmtMatchSep}{0.28cm}
\newcommand{\matchHead}[1]{\par\noindent\textit{#1}\par\nopagebreak\vspace{0.2cm}}
% шаблонний заголовок стовпця зі старого макета -- лише коли завдання не має власного \matchHead
\makeatletter\newcommand{\nmtHeadUnless}[2]{\in@{\matchHead}{#1}\ifin@\else#2\fi}\makeatother
\newcommand{\matchItem}[2]{\par\noindent\makebox[\nmtMatchLab][l]{\textbf{#1}}%
\parbox[t]{\dimexpr\linewidth-\nmtMatchLab\relax}{\raggedright #2}\par\nopagebreak\vspace{\nmtMatchSep}}
\newcommand{\ansTheme}[1]{\par\vspace{0.25cm}\noindent{\bfseries\color{mainGreen}#1}\par\vspace{0.12cm}}
\newcommand{\ansType}[1]{\par\vspace{0.1cm}\noindent{\itshape\color{yearOrange}#1}\par\vspace{0.08cm}}

\graphicspath{{../1. Числа. Звичайні та десяткові дроби. Модуль числа/}{../10. Основи планіметрії/}{../11. Трикутники/}{../12. Паралелограм/}{../13. Прямокутник/}{../14. Квадрат/}{../15. Ромб/}{../16. Трапеція/}{../17. Коло, круг і їх елементи/}{../18. Координати і вектори на площині і в просторі. Рівняння кола/}{../19. Лінійні та квадратні нерівності. Системи лінійних нерівностей. Метод інтервалів/}{../2.відсотки, арифметичні задачі, подільність/}{../20. Функція та її властивості/}{../21.  Графіки елементарних функцій, їх побудова графіків функцій та їх перетворення/}{../22. Модульні  рівняння та нерівності/}{../23. Ірраціональні рівняння та нерівності/}{../24. Арифметична прогресія/}{../25. Геометрична прогресія/}{../26. Тригонометричні вирази/}{../27. Тригонометричні рівняння/}{../28. Показникова функція. Показникові рівняння/Показникові рівняння/}{../28. Показникова функція. Показникові рівняння/Показникові функція і вирази/}{../29. Показникові нерівності/}{../3. Степені. Одночлени. стандартний вигляд числа/}{../30. Логарифм. Логарифмічна функція/}{../31. Логарифмічні рівняння/}{../32. Логарифмічні нерівності/}{../33. Похідна/}{../34. Первісна/}{../35. Комбінаторика/}{../36. Теорія ймовірності/}{../37. Статистика/}{../38. Призма. Паралелепіпед. Куб/}{../39. Піраміда/}{../4.Многочлени. ФСМ/}{../40. Циліндр/}{../41. Конус/}{../42. Куля. Сфера/}{../43. Параметри/}{../5. дробові раціональні вирази і перетворення/}{../6. Лінійні рівняння/}{../7. Квадратні корені. Корені вищих степенів. Раціональний степінь/}{../8. квадратні рівняння, і рівняння, що до них зводяться/}{../9. Системи рівнянь/}}
\renewcommand{\instructionBox}[1]{%
\par\vspace{0.3cm}\noindent
\begin{tcolorbox}[nobeforeafter, width=\linewidth, colback=instrBg, colframe=instrBorder, boxrule=0.6pt, arc=2pt,
    left=8pt, right=8pt, top=4pt, bottom=4pt]
\centering\small\bfseries #1
\end{tcolorbox}
\par\nopagebreak\vspace{0.15cm}\nopagebreak
}
\renewcommand{\nmtAnswerBox}{%
\par\nopagebreak\vspace{0.2cm}\noindent
Відповідь:\ \
\begin{tabular}{@{}*{5}{|>{\centering\arraybackslash}p{0.8cm}}|@{\hspace{0.1cm}\textbf{,}\hspace{0.1cm}}*{3}{|>{\centering\arraybackslash}p{0.8cm}}|@{}}
\hline
\rule[-0.2cm]{0pt}{0.7cm} & & & & & & & \\
\hline
\end{tabular}
\par\vspace{0.3cm}
}

\begin{document}
\chapterTitle{Пробний НМТ з математики № 1: відповіді}

\begin{center}\renewcommand{\arraystretch}{1.25}
\begin{tabular}{|c|c||c|c||c|c|}\hline
\textbf{№} & \textbf{Відповідь} & \textbf{№} & \textbf{Відповідь} & \textbf{№} & \textbf{Відповідь} \\ \hline

1 & В & 9 & В & 17 & 1~--~А, 2~--~Г, 3~--~В \\ \hline
2 & Г & 10 & А & 18 & 1~--~В, 2~--~Д, 3~--~А \\ \hline
3 & Д & 11 & Д & 19 & $0{{,}}8$ \\ \hline
4 & А & 12 & А & 20 & $-5{{,}}2$ \\ \hline
5 & Б & 13 & Б & 21 & $480$ \\ \hline
6 & Г & 14 & Г & 22 & $-18$ \\ \hline
7 & Б & 15 & Д &  &  \\ \hline
8 & В & 16 & 1~--~В, 2~--~А, 3~--~Б &  &  \\ \hline
\end{tabular}\end{center}

\noindent{\small\color{gray!80!black}Оцінювання: 1--15 -- по 1 балу; 16--18 -- по 1 балу за кожну правильно встановлену відповідність (до 3 балів); 19--22 -- по 2 бали. Разом 32 бали.}\par\vspace{0.3cm}

\sectionTitle{Розв'язки й типові помилки}

\par\noindent\textbf{1.} $7=3+3+1$: платять за $2+2+1=5$ йогуртів, $5\cdot36=180$~грн.\quad{\small\color{gray!80!black}(Текстова задача: найменша сума (набори й поштучно))}\par\nopagebreak

\noindent{\footnotesize\color{gray!85!black}Типові помилки: А -- за сьомий йогурт теж не заплатили; Б -- знижку $\frac13$ застосували до всіх семи йогуртів; Г -- безкоштовним вважали лише один йогурт; Д -- не врахували акції.}\par\vspace{0.15cm}

\par\noindent\textbf{2.} У суботу продали $260$ квитків, у середу -- $100$; різниця $160$.\quad{\small\color{gray!80!black}(Стовпчикова діаграма)}\par\nopagebreak

\noindent{\footnotesize\color{gray!85!black}Типові помилки: А -- відняли найменше значення (вівторок); Б -- сплутали суботу з неділею; В -- порівняли суботу з п'ятницею (сусідній стовпчик); Д -- виписали кількість квитків у суботу.}\par\vspace{0.15cm}

\par\noindent\textbf{3.} $BE\parallel AK$, $AB$ -- січна, тому $\angle ABE=180^\circ-70^\circ=110^\circ$ (односторонні кути). $\alpha=\angle ABE+\angle EBC=110^\circ+30^\circ=140^\circ$.\quad{\small\color{gray!80!black}(Практична задача: настільна лампа (паралельні прямі й січна))}\par\nopagebreak

\noindent{\footnotesize\color{gray!85!black}Типові помилки: А -- відняли кути: $70^\circ-30^\circ$; Б -- від $\angle ABE=110^\circ$ відняли $30^\circ$, а не додали (ніби промінь $BC$ лежить між $BA$ і $BE$): $110^\circ-30^\circ$; В -- сплутали односторонні кути з різносторонніми ($\angle ABE=70^\circ$): $70^\circ+30^\circ$; Г -- знайшли лише кут $ABE$, не додавши $30^\circ$.}\par\vspace{0.15cm}

\par\noindent\textbf{4.} $16=0{,}5^{-4}$; основа менша за $1$, тому $x^2-4x\leqslant-4$, тобто $(x-2)^2\leqslant0$ -- лише $x=2$.\quad{\small\color{gray!80!black}(Показникова: зведення до основи (основа менша за 1))}\par\nopagebreak

\noindent{\footnotesize\color{gray!85!black}Типові помилки: Б -- розв'язали $(x-2)^2\leqslant0$ як $x-2\leqslant0$; В -- розв'язали $(x-2)^2\leqslant0$ як $x-2\geqslant0$; Г -- забули про рівність: квадрат не буває від'ємним; Д -- не змінили знак нерівності (основа $0{,}5<1$): $(x-2)^2\geqslant0$.}\par\vspace{0.15cm}

\par\noindent\textbf{5.} Пряма паралельна площині, якщо вона паралельна якійсь прямій цієї площини й не лежить у ній. $AB\parallel A_1B_1$ -- верхня основа; $AB\parallel DE$ (протилежні сторони правильного шестикутника) -- грань $DEE_1D_1$. Нижня основа й грань $ABB_1A_1$ містять $AB$, грані $BCC_1B_1$ і $AFF_1A_1$ мають з нею спільну точку, а прямі $AB$ і $CD$ (так само $AB$ і $EF$) лежать в одній площині основи й не паралельні, тож перетинаються -- отже, $AB$ перетинає площини $CDD_1C_1$ і $EFF_1E_1$. Отже, таких граней дві.\quad{\small\color{gray!80!black}(Куб і призма: пряма, паралельна площині грані)}\par\nopagebreak

\noindent{\footnotesize\color{gray!85!black}Типові помилки: А -- урахували лише верхню основу, пропустили бічну грань $DEE_1D_1$ ($DE\parallel AB$); В -- порахували ребра, паралельні $AB$ ($A_1B_1$, $DE$, $D_1E_1$), а не грані; Г -- зарахували грані $CDD_1C_1$ і $EFF_1E_1$, бо відрізок $AB$ їх не торкається (або зарахували грані, що містять $AB$: нижню основу й $ABB_1A_1$); Д -- вважали паралельними всі грані, що не містять ребра $AB$.}\par\vspace{0.15cm}

\par\noindent\textbf{6.} $(x+3)^2-(x+1)(x+5)=x^2+6x+9-(x^2+6x+5)=4$ -- не залежить від $x$. Решта: $12x$; $6x-9$; $3x+9$; $9-12x$.\quad{\small\color{gray!80!black}(Формули скороченого множення: спрощення)}\par\nopagebreak

\noindent{\footnotesize\color{gray!85!black}Типові помилки: А -- загубили подвоєний добуток: $(x^2+9)-(x^2+9)=0$; насправді $12x$; Б -- прийняли $(x-3)(3-x)$ за різницю квадратів $9-x^2$: $9-x^2+x^2=9$; насправді $6x-9$; В -- не подвоїли добуток: $(x+3)^2=x^2+3x+9$, $x^2+3x+9-x^2-3x=9$; насправді $3x+9$; Д -- мінус перед дужками застосували лише до $x^2$: $x^2-6x+9-x^2+6x=9$; насправді $9-12x$.}\par\vspace{0.15cm}

\par\noindent\textbf{7.} Заміна $t=x^2\geqslant0$: $t^2-t-72=0$, $t=9$ або $t=-8$ (не підходить). $x^2=9$, $x=\pm3$, добуток коренів $3\cdot(-3)=-9$.\quad{\small\color{gray!80!black}(Біквадратне рівняння: добуток коренів)}\par\nopagebreak

\noindent{\footnotesize\color{gray!85!black}Типові помилки: А -- застосували теорему Вієта до $x^4-x^2-72=0$, як до квадратного: добуток $t_1t_2=-72$; В -- знайшли суму коренів замість добутку; Г -- узяли корінь $t=9$ допоміжного рівняння замість добутку $x_1x_2=3\cdot(-3)$; Д -- не відкинули $x^2=-8$ і записали ще корені $\pm\sqrt8$: $3\cdot(-3)\cdot\sqrt8\cdot(-\sqrt8)=72$.}\par\vspace{0.15cm}

\par\noindent\textbf{8.} Нулі функції $y=f(x)-2$ -- це корені рівняння $f(x)=2$, тобто абсциси спільних точок графіка з прямою $y=2$. Ця пряма дотикається до графіка в точці $(-4;\,2)$ і перетинає його ще у двох точках (між $-2$ і $0$ та між $0$ і $3$). Отже, нулів три.\quad{\small\color{gray!80!black}(Нулі функції за графіком)}\par\nopagebreak

\noindent{\footnotesize\color{gray!85!black}Типові помилки: А -- розв'язали $f(x)=-2$: графік перенесли вгору, а не вниз; Б -- не врахували точку $(-4;\,2)$, у якій пряма $y=2$ лише дотикається до графіка; Г -- точку дотику $(-4;\,2)$ порахували як дві точки перетину (по одній з кожного боку вершини); Д -- порахували нулі самої функції $y=f(x)$, не врахувавши $-2$.}\par\vspace{0.15cm}

\par\noindent\textbf{9.} $6^2+8^2=10^2$ -- прямокутний; $R=\frac{10}{2}=5$; площа $\frac12\cdot6\cdot8=24\ne48$. Правильні I і II.\quad{\small\color{gray!80!black}(Трикутник: медіана, бісектриса, висота, нерівність трикутника)}\par\nopagebreak

\noindent{\footnotesize\color{gray!85!black}Типові помилки: А -- не знали, що центр описаного кола прямокутного трикутника -- середина гіпотенузи; Б -- перевірили лише площу, рахуючи її без $\frac12$; Г -- вважали трикутник не прямокутним, а площу -- $6\cdot8$; Д -- площу рахували як $6\cdot8$ без $\frac12$.}\par\vspace{0.15cm}

\par\noindent\textbf{10.} $b_5=b_2\cdot q^3$, тож $q^3=\frac{-2}{16}=-\frac18$, $q=-\frac12$; $b_1=\frac{b_2}{q}=16:\left(-\frac12\right)=-32$. Перші чотири члени: $-32;\ 16;\ -8;\ 4$, їхня сума $-20$ (або за формулою $S_4=\frac{b_1(q^4-1)}{q-1}=\frac{-32\cdot\left(\frac1{16}-1\right)}{-\frac32}=-20$).\quad{\small\color{gray!80!black}(Геометрична прогресія: знаменник з $b_2$ і $b_5$, сума перших членів)}\par\nopagebreak

\noindent{\footnotesize\color{gray!85!black}Типові помилки: Б -- перший член знайшли як $b_2\cdot q=-8$ (помножили на знаменник замість поділити); В -- суму почали з $b_2$: знайшли $b_2+b_3+b_4+b_5=16-8+4-2$; Г -- знаменник знайшли з відношення $\frac{b_2}{b_5}$: $q^3=-8$, $q=-2$, тоді $b_1=-8$; Д -- загубили знак: $\sqrt[3]{-\frac18}=\frac12$, тоді $b_1=32$ і всі члени додатні.}\par\vspace{0.15cm}

\par\noindent\textbf{11.} $P=\frac{5}{20}=\frac14$.\quad{\small\color{gray!80!black}(Класична ймовірність; протилежна подія)}\par\nopagebreak

\noindent{\footnotesize\color{gray!85!black}Типові помилки: А -- знайшли ймовірність протилежної події; Б -- поділили $5$ на кількість \textit{інших} треків ($15$); В -- «один трек із двадцяти»; Г -- вважали, що треків $5+20=25$.}\par\vspace{0.15cm}

\par\noindent\textbf{12.} $B(4;\,0;\,0)$, $D(0;\,8;\,0)$, $D_1(0;\,8;\,6)$, тож $M\left(0;\,8;\,\frac{0+6}{2}\right)=M(0;\,8;\,3)$. $\overrightarrow{BM}=(0-4;\,8-0;\,3-0)=(-4;\,8;\,3)$.\quad{\small\color{gray!80!black}(Многогранник у системі координат: координати вершин і векторів)}\par\nopagebreak

\noindent{\footnotesize\color{gray!85!black}Типові помилки: Б -- знайшли координати вектора $\overrightarrow{MB}$ (від координат $B$ відняли координати $M$); В -- записали координати точки $M$ (не відняли координати початку вектора $B$); Г -- узяли вершину $D_1$ замість середини ребра $DD_1$; Д -- координати середини ребра взяли як половини координат $D_1$: $M(0;\,4;\,3)$.}\par\vspace{0.15cm}

\par\noindent\textbf{13.} $\sin^2\alpha=1-0{,}3=0{,}7$; $5-4\cdot0{,}7=2{,}2$.\quad{\small\color{gray!80!black}(Основна тотожність: значення за відомим квадратом)}\par\nopagebreak

\noindent{\footnotesize\color{gray!85!black}Типові помилки: А -- знайшли $4\sin^2\alpha$, не віднявши від $5$; В -- підставили $\cos^2\alpha$ замість $\sin^2\alpha$; Г -- знайшли лише $\sin^2\alpha$; Д -- «$5-4\sin^2\alpha=1-4\cos^2\alpha$».}\par\vspace{0.15cm}

\par\noindent\textbf{14.} Висота трикутника $AKD$, проведена з $K$, дорівнює $AB=4\sqrt3$. З трикутника з кутом $30^\circ$ при $D$: $KD=\frac{4\sqrt3}{\sin30^\circ}=8\sqrt3$; з прямокутного трикутника $AKD$: $AD=\frac{KD}{\cos30^\circ}=16$. Площа $\frac12\cdot16\cdot4\sqrt3=32\sqrt3$~см$^2$.\quad{\small\color{gray!80!black}(Складна обчислювальна задача (кілька кроків, кути $30^\circ$ і $60^\circ$))}\par\nopagebreak

\noindent{\footnotesize\color{gray!85!black}Типові помилки: А -- $AD$ знайшли як $AB:\sin60^\circ=8$; Б -- знайшли площу прямокутника; В -- $AD$ знайшли як $KD\cdot\cos30^\circ=12$ (помножили замість поділити); Д -- вважали, що $AD=2AB$ (катет проти кута $30^\circ$).}\par\vspace{0.15cm}

\par\noindent\textbf{15.} Перемножимо рівняння: $x^3y^3=12\cdot18=216$, тобто $(xy)^3=216$, звідки $x_0y_0=\sqrt[3]{216}=6$. Перевірка: $x^2y=x\cdot xy$, тож $x_0=\frac{12}{6}=2$, $y_0=\frac{18}{6}=3$.\quad{\small\color{gray!80!black}(Нелінійна система: перемноження рівнянь, $(xy)^3=216$)}\par\nopagebreak

\noindent{\footnotesize\color{gray!85!black}Типові помилки: А -- перемножили рівняння, але добули квадратний корінь замість кубічного: $\sqrt{216}$; Б -- $\sqrt[3]{216}$ знайшли як $216:3$; В -- перемножили рівняння й не добули корінь: знайшли $(x_0y_0)^3$; Г -- поділили друге рівняння на перше й зупинилися: знайшли $y_0:x_0$.}\par\vspace{0.15cm}

\par\noindent\textbf{16.} Ключ: $x+2=\sqrt3$, а $x<0$, бо $\sqrt3<2$. 1) $(x+2)^{-2}=(\sqrt3)^{-2}=\frac13$; 2) $\frac{1}{\sqrt3-2}=\frac{\sqrt3+2}{(\sqrt3-2)(\sqrt3+2)}=\frac{\sqrt3+2}{3-4}=-2-\sqrt3$; 3) $(x+2)^2=3$, тож $\log_3\frac{1}{(x+2)^2}=\log_3\frac13=\log_3 3^{-1}=-1$.\quad{\small\color{gray!80!black}(Значення виразу за заданого значення змінної)}\par\nopagebreak

\noindent{\footnotesize\color{gray!85!black}Типові помилки: 1~--~Д: не врахували мінус у показнику: $(\sqrt3)^2=3$; 2~--~Б: помножили на спряжений вираз лише знаменник: $\frac{1}{(\sqrt3-2)(\sqrt3+2)}=\frac{1}{3-4}$; 3~--~В: записали число під логарифмом замість значення логарифма: $\frac{1}{(\sqrt3)^2}=\frac13$; 3~--~Г: загубили дріб: $\log_3(\sqrt3)^{2}=\log_3 3=1$ замість $\log_3\frac13$.}\par\vspace{0.15cm}

\par\noindent\textbf{17.} 1) парна: $f(-x)=f(x)$; 2) під логарифмом -- лише додатні числа: $x-2>0$, тож $D(y)=(2;\,+\infty)$, а $E(y)$ -- усі числа; 3) кутовий коефіцієнт $-2<0$, функція спадає.\quad{\small\color{gray!80!black}(Функція -- її властивість: парність, монотонність, період, область визначення)}\par\nopagebreak

\noindent{\footnotesize\color{gray!85!black}Типові помилки: 1~--~В: сплутали: парабола спадає лише на $(-\infty;\,0]$; 2~--~Д: сплутали область визначення з областю значень; 3~--~Б: вважали лінійну функцію непарною.}\par\vspace{0.15cm}

\par\noindent\textbf{18.} $AD=\frac{5}{3}\cdot12=20$~см. $BC=MN=12$, $KD=BC=12$ (протилежні сторони паралелограма), тож $AK=20-12=8$~см -- сторона трикутника, і $AM=8$~см. 1) $S_{ABK}=\frac{8^2\sqrt3}{4}=16\sqrt3$~см$^2$. 2) $\angle BKD=180^\circ-60^\circ=120^\circ$, $S_{KBCD}=8\cdot12\cdot\sin120^\circ=48\sqrt3$~см$^2$ (або висота $4\sqrt3$ до сторони $KD=12$). 3) $S_{ADNM}=\frac{20+12}{2}\cdot8=128$~см$^2$.\quad{\small\color{gray!80!black}(Площі фігур на рисунку: рівносторонній трикутник, паралелограм, прямокутна трапеція)}\par\nopagebreak

\noindent{\footnotesize\color{gray!85!black}Типові помилки: 1~--~Г: площу трикутника рахували як $a\cdot h=8\cdot4\sqrt3$, без $\frac12$; 2~--~Г: висоту $4\sqrt3$ помножили на сторону $KB=8$, а не на сторону $KD=12$, до якої її проведено; 3~--~Б: за висоту трапеції взяли $MN=12$ замість $AM=8$: $\frac{20+12}{2}\cdot12$.}\par\vspace{0.15cm}

\par\noindent\textbf{19.} Пряма пропорційність $g(x)=kx$; з $6=k\cdot(-3)$ маємо $k=-2$, тож $g(x)=-2x$ і $g'(x)=-2$ для всіх $x$, зокрема $g'(-1)=-2$. $f(2)=2$ (бо $2>0$), $f(-1)=5$ (бо $-1<0$). Тоді $\dfrac{-2+3\cdot2}{5}=\dfrac45=0{,}8$.\quad{\small\color{gray!80!black}(Кусково-стала функція й пряма пропорційність: значення виразу з $g'(-1)$)}\par\nopagebreak

\noindent{\footnotesize\color{gray!85!black}Типові помилки: $1{,}6$ -- замість похідної $g'(-1)=-2$ підставили значення функції $g(-1)=2$; $1{,}1$ -- кутовий коефіцієнт знайшли як $\frac{x}{y}=-0{,}5$; $6{,}5$ -- переплутали значення $f$ для додатних і від'ємних $x$; $1{,}2$ -- спершу підставили $x=-1$, а потім узяли похідну від числа: $\big(g(-1)\big)'=0$.}\par\vspace{0.15cm}

\par\noindent\textbf{20.} Сума температур шести днів $-2-6-3-7-5-4=-27$, середнє $\frac{-27}{6}=-4{,}5$, середнє за тиждень $-4{,}5-0{,}1=-4{,}6$. Сума за сім днів $7\cdot(-4{,}6)=-32{,}2$, тож у неділю $-32{,}2-(-27)=-5{,}2^\circ$C. (Інакше: щоб середнє всіх семи днів стало на $0{,}1^\circ$ нижчим за $-4{,}5$, неділя має бути нижчою за $-4{,}5$ на $7\cdot0{,}1=0{,}7^\circ$: $-4{,}5-0{,}7=-5{,}2$.)\quad{\small\color{gray!80!black}(Нове значення знизило середнє на відому величину: знайти це значення (від'ємні дробові))}\par\nopagebreak

\noindent{\footnotesize\color{gray!85!black}Типові помилки: $-4{,}6$ -- вважали, що температура в неділю дорівнює новому середньому $-4{,}6$; $-5{,}1$ -- знизили на $0{,}1^\circ$ лише шість попередніх днів, не врахувавши, що й неділя входить у нове середнє; $-5{,}3$ -- у способі «середнє $-\,7\cdot0{,}1$» узяли нове середнє $-4{,}6$ замість старого $-4{,}5$; $-4{,}7$ -- усереднили старе середнє й температуру неділі, ніби це дві рівноправні величини: $\frac{-4{,}5+y}{2}=-4{,}6$; $-3{,}8$ -- помилилися в знаку: середнє «нижче», а до старого середнього додали $7\cdot0{,}1$.}\par\vspace{0.15cm}

\par\noindent\textbf{21.} Відстань від центра основи піраміди до її сторони $4$~см, тому апофема $\frac{4}{\cos60^\circ}=8$~см, а бічна поверхня піраміди $4\cdot\frac12\cdot8\cdot8=128$~см$^2$. Бічна поверхня призми $4\cdot8\cdot9=288$~см$^2$, нижня основа $64$~см$^2$ (верхня основа призми всередині тіла). Разом $288+64+128=480$~см$^2$.\quad{\small\color{gray!80!black}(Призма й піраміда)}\par\nopagebreak

\noindent{\footnotesize\color{gray!85!black}Типові помилки: $544$ -- до поверхні тіла додали й верхню основу призми (вона всередині тіла); $416$ -- забули нижню основу призми; $384$ -- апофему знайшли як $4\cdot\cos60^\circ=2$; $608$ -- додали повні поверхні призми й піраміди (спільну грань порахували двічі) або площу грані піраміди рахували без $\frac12$.}\par\vspace{0.15cm}

\par\noindent\textbf{22.} Нехай $t=2^{x-a}>0$: $t^2-2t-8=0$, $t=4$ ($t=-2$ не підходить), тож $x=a+2$. Знаменник визначений і не нуль при $-5<x<5$: $-7<a<3$, цілі $a=-6,\ldots,2$, сума $-18$.\quad{\small\color{gray!80!black}(Дріб дорівнює нулю: корінь чисельника має потрапити в ОДЗ знаменника)}\par\nopagebreak

\noindent{\footnotesize\color{gray!85!black}Типові помилки: $-22$ -- під коренем у знаменнику взяли нестрогу нерівність; $18$ -- помилилися в знаку: $x=a-2$; $-36$ -- з $2^{x-a}=4$ отримали $x-a=4$.}\par\vspace{0.15cm}

\end{document}
